\chapter{2017年天津卷（文）}

\section{单选题}

\begin{problem}
设集合 \(A=\{1,2,6\}\)，\(B=\{2,4\}\)，\(C=\{1,2,3,4\}\)，则 \((A\cup B)\cap C=\)（\quad）

\choices
    {\(\{2\}\)}
    {\(\{1,2,4\}\)}
    {\(\{1,2,4,6\}\)}
    {\(\{1,2,3,4,6\}\)}
\end{problem}

\begin{answer}
B
\end{answer}

\begin{solution}
先求并集，\(A\cup B=\{1,2,4,6\}\).再取交集，\((A\cup B)\cap C=\{1,2,4\}\)，因此选 B.
\end{solution}

\begin{problem}
设 \(x\in\R\)，则“\(2-x\ge 0\)”是“\(|x-1|\le 1\)”的（\quad）

\choices
    {充分而不必要条件}
    {必要而不充分条件}
    {充要条件}
    {既不充分也不必要条件}
\end{problem}

\begin{answer}
B
\end{answer}

\begin{solution}
由 \(2-x\ge0\) 得 \(x\le2\).由 \(|x-1|\le1\) 得 \(0\le x\le2\)，所以 \(|x-1|\le1\) 一定推出 \(2-x\ge0\)，即前者是后者的必要条件.取 \(x=-1\)，有 \(2-x=3\ge0\)，但 \(|x-1|=2>1\)，故它不是充分条件.因此选 B.
\end{solution}

\begin{problem}
有 \(5\) 支彩笔（除颜色外无差别），颜色分别为红，黄，蓝，绿，紫. 从这 \(5\) 支彩笔中任取 \(2\) 支不同颜色的彩笔，则取出的 \(2\) 支彩笔中含有红色彩笔的概率为（\quad）

\choices
    {\(\dfrac45\)}
    {\(\dfrac35\)}
    {\(\dfrac25\)}
    {\(\dfrac15\)}
\end{problem}

\begin{answer}
C
\end{answer}

\begin{solution}
从五支彩笔中任取两支的基本事件数为 \(\mathrm C_5^2=10\).若取出的两支中含红色彩笔，另一支可从其余四支中任选一支，故有 \(4\) 个有利事件.因此所求概率为 \(\dfrac{4}{10}=\dfrac25\)，选 C.
\end{solution}

\begin{problem}
阅读如图的程序框图，运行相应的程序，若输入 \(N\) 的值为 \(19\)，则输出 \(N\) 的值为（\quad）

\bitmapfigure[max width=5.2cm]{img/2017/tianjin_liberal/q4_fig1.png}

\choices
    {\(0\)}
    {\(1\)}
    {\(2\)}
    {\(3\)}
\end{problem}

\begin{answer}
C
\end{answer}

\begin{solution}
按框图逐次执行.初始 \(N=19\)，\(19\) 不能被 \(3\) 整除，执行 \(N=N-1\)，得 \(N=18\)；\(18\) 能被 \(3\) 整除，执行 \(N=N/3\)，得 \(N=6\)；\(6\) 能被 \(3\) 整除，执行 \(N=N/3\)，得 \(N=2\).此时 \(N\le3\)，输出 \(2\)，故选 C.
\end{solution}

\begin{problem}
已知双曲线 \(\dfrac{x^2}{a^2}-\dfrac{y^2}{b^2}=1\)（\(a>0\)，\(b>0\)）的右焦点为 \(F\)，点 \(A\) 在双曲线的渐近线上，\(\triangle OAF\) 是边长为 \(2\) 的等边三角形（\(O\) 为原点），则双曲线的方程为（\quad）

\choices
    {\(\dfrac{x^2}{4}-\dfrac{y^2}{12}=1\)}
    {\(\dfrac{x^2}{12}-\dfrac{y^2}{4}=1\)}
    {\(\dfrac{x^2}{3}-y^2=1\)}
    {\(x^2-\dfrac{y^2}{3}=1\)}
\end{problem}

\begin{answer}
D
\end{answer}

\begin{solution}
设双曲线的半焦距为 \(c\)，则右焦点为 \(F(c,0)\).由于 \(\triangle OAF\) 是边长为 \(2\) 的等边三角形，所以 \(c=OF=2\)，且 \(OA=2\).因此点 \(A\) 可取 \((1,\sqrt3)\) 或 \((1,-\sqrt3)\)，从而渐近线的斜率的绝对值为 \(\sqrt3\).双曲线的渐近线方程为 \(y=\pm\dfrac ba x\)，故 \(\dfrac ba=\sqrt3\).

又因为 \(c^2=a^2+b^2\)，所以 \(4=a^2+3a^2\)，得 \(a=1\)，\(b=\sqrt3\).故双曲线方程为 \(x^2-\dfrac{y^2}{3}=1\)，选 D.
\end{solution}

\begin{problem}
已知奇函数 \(f(x)\) 在 \(\R\) 上是增函数. 若 \(a=-f\!\left(\log_2\dfrac15\right)\)，\(b=f(\log_2 4.1)\)，\(c=f(2^{0.8})\)，则 \(a\)，\(b\)，\(c\) 的大小关系为（\quad）

\choices
    {\(a<b<c\)}
    {\(b<a<c\)}
    {\(c<b<a\)}
    {\(c<a<b\)}
\end{problem}

\begin{answer}
C
\end{answer}

\begin{solution}
由 \(f\) 为奇函数，\(\log_2\dfrac15=-\log_2 5\)，所以
\(a=-f(-\log_2 5)=f(\log_2 5)\).又因为 \(4<4.1<5\)，故 \(\log_2 4.1<\log_2 5\)，从而 \(b<a\).

因为 \(0.8<1\)，所以 \(2^{0.8}<2\)；又因为 \(4.1>4=2^2\)，所以 \(\log_2 4.1>2>2^{0.8}\).由 \(f\) 在 \(\R\) 上递增，得 \(c<b\).因此 \(c<b<a\)，选 C.
\end{solution}

\begin{problem}
设函数 \(f(x)=2\sin(\omega x+\varphi)\)，\(x\in\R\)，其中 \(\omega>0\)，\(|\varphi|<\pi\). 若 \(f\!\left(\dfrac{5\pi}{8}\right)=2\)，\(f\!\left(\dfrac{11\pi}{8}\right)=0\)，且 \(f(x)\) 的最小正周期大于 \(2\pi\)，则（\quad）

\choices
    {\(\omega=\dfrac23\)，\(\varphi=\dfrac\pi{12}\)}
    {\(\omega=\dfrac23\)，\(\varphi=-\dfrac{11\pi}{12}\)}
    {\(\omega=\dfrac13\)，\(\varphi=-\dfrac{11\pi}{24}\)}
    {\(\omega=\dfrac13\)，\(\varphi=\dfrac{7\pi}{24}\)}
\end{problem}

\begin{answer}
A
\end{answer}

\begin{solution}
由最小正周期 \(\dfrac{2\pi}{\omega}>2\pi\)，得 \(0<\omega<1\).设整数 \(k\)，\(m\)，则
\(\dfrac{5\pi}{8}\omega+\varphi=\dfrac\pi2+2k\pi\)，\(\dfrac{11\pi}{8}\omega+\varphi=m\pi\).两式相减得
\(\dfrac{3\pi}{4}\omega=\left(m-2k-\dfrac12\right)\pi\).由于 \(0<\omega<1\)，只能取 \(m-2k=1\)，从而 \(\omega=\dfrac23\).代回第一式，得 \(\varphi=\dfrac\pi{12}+2k\pi\).结合 \(|\varphi|<\pi\)，得 \(\varphi=\dfrac\pi{12}\).故选 A.
\end{solution}

\begin{problem}
已知函数 \(f(x)\) 满足：当 \(x<1\) 时，\(f(x)=|x|+2\)；当 \(x\ge 1\) 时，\(f(x)=x+\dfrac2x\). 设 \(a\in\R\)，若关于 \(x\) 的不等式 \(f(x)\ge \left|\dfrac x2+a\right|\) 在 \(\R\) 上恒成立，则 \(a\) 的取值范围是（\quad）

\choices
    {\([-2,2]\)}
    {\([-2\sqrt3,2]\)}
    {\([-2,2\sqrt3]\)}
    {\([-2\sqrt3,2\sqrt3]\)}
\end{problem}

\begin{answer}
A
\end{answer}

\begin{solution}
当 \(x<0\) 时，\(f(x)=2-x>0\)，故不等式等价于
\(-2+x\le\dfrac x2+a\le2-x\).于是对所有 \(x<0\)，有 \(a\ge\dfrac x2-2\) 且 \(a\le2-\dfrac{3x}{2}\).令 \(x\) 从左侧趋近于 \(0\)，得到必要条件 \(a\ge-2\) 且 \(a\le2\).

当 \(0\le x<1\) 时，\(f(x)=x+2\)，不等式等价于 \(-x-2\le\dfrac x2+a\le x+2\)，因而 \(a\ge-2-\dfrac{3x}{2}\) 且 \(a\le2+\dfrac x2\)，对 \(0\le x<1\) 仍只给出 \(-2\le a\le2\).

当 \(x\ge1\) 时，\(f(x)=x+\dfrac2x>0\)，不等式等价于
\(-x-\dfrac2x\le\dfrac x2+a\le x+\dfrac2x\).因此
\(a\ge-\dfrac{3x}{2}-\dfrac2x\) 且 \(a\le\dfrac x2+\dfrac2x\).其中 \(-\dfrac{3x}{2}-\dfrac2x\le-2\sqrt3<-2\)，而由基本不等式 \(\dfrac x2+\dfrac2x\ge2\)，故这部分不会比 \(-2\le a\le2\) 产生更强的限制.

所以当且仅当 \(a\in[-2,2]\) 时原不等式对所有实数 \(x\) 恒成立，选 A.
\end{solution}

\section{填空题}

\begin{problem}
\(a\in\R\)，\(\i\) 为虚数单位，若 \(\dfrac{a-\i}{2+\i}\) 为实数，则 \(a\) 的值为\fillinblank{}.
\end{problem}

\begin{answer}
\(-2\)
\end{answer}

\begin{solution}
分母乘以共轭复数，得
\(\dfrac{a-\i}{2+\i}=\dfrac{(a-\i)(2-\i)}{5}=\dfrac{2a-1-(a+2)\i}{5}\).该复数为实数，当且仅当虚部为零，即 \(a+2=0\)，所以 \(a=-2\).
\end{solution}

\begin{problem}
已知 \(a\in\R\)，设函数 \(f(x)=ax-\ln x\) 的图象在点 \((1,f(1))\) 处的切线为 \(l\)，则 \(l\) 在 \(y\) 轴上的截距为\fillinblank{}.
\end{problem}

\begin{answer}
\(1\)
\end{answer}

\begin{solution}
函数定义域为 \(x>0\)，且 \(f(1)=a\)，\(f'(x)=a-\dfrac1x\)，所以切线斜率为 \(f'(1)=a-1\).切线方程为 \(y-a=(a-1)(x-1)\)，化简得 \(y=(a-1)x+1\).令 \(x=0\)，得到在 \(y\) 轴上的截距为 \(1\).
\end{solution}

\begin{problem}
已知一个正方体的所有顶点在一个球面上，若这个正方体的表面积为 \(18\)，则这个球的体积为\fillinblank{}.
\end{problem}

\begin{answer}
\(\dfrac{9\pi}{2}\)
\end{answer}

\begin{solution}
设正方体棱长为 \(s\)，则 \(6s^2=18\)，得 \(s=\sqrt3\).球为正方体的外接球，其直径等于正方体的体对角线 \(s\sqrt3=3\)，故半径 \(r=\dfrac32\).所以球的体积为 \(\dfrac43\pi r^3=\dfrac43\pi\left(\dfrac32\right)^3=\dfrac{9\pi}{2}\).
\end{solution}

\begin{problem}
设抛物线 \(y^2=4x\) 的焦点为 \(F\)，准线为 \(l\). 已知点 \(C\) 在 \(l\) 上，以 \(C\) 为圆心的圆与 \(y\) 轴的正半轴相切于点 \(A\). 若 \(\angle FAC=120^\circ\)，则圆的方程为\fillinblank{}.
\end{problem}

\begin{answer}
\(\left(x+1\right)^2+\left(y-\sqrt{3}\right)^2=1\)
\end{answer}

\begin{solution}
由 \(y^2=4x\) 得焦点 \(F(1,0)\)，准线 \(l\) 的方程为 \(x=-1\).设 \(C(-1,t)\).圆以 \(C\) 为圆心且与 \(y\) 轴相切，所以半径为 \(1\)，切点为 \(A(0,t)\).由于切点在 \(y\) 轴正半轴上，故 \(t>0\).

在点 \(A\) 处，\(\overrightarrow{AF}=(1,-t)\)，\(\overrightarrow{AC}=(-1,0)\)，所以
\(\cos\angle FAC=\dfrac{\overrightarrow{AF}\cdot\overrightarrow{AC}}{|AF||AC|}=-\dfrac1{\sqrt{1+t^2}}\).由 \(\angle FAC=120^\circ\)，得 \(-\dfrac1{\sqrt{1+t^2}}=-\dfrac12\)，从而 \(t^2=3\).结合 \(t>0\)，得 \(t=\sqrt3\).因此圆的方程为 \(\left(x+1\right)^2+\left(y-\sqrt3\right)^2=1\).
\end{solution}

\begin{problem}
若 \(a\)，\(b\in\R\)，\(ab>0\)，则 \(\dfrac{a^4+4b^4+1}{ab}\) 的最小值为\fillinblank{}.
\end{problem}

\begin{answer}
\(4\)
\end{answer}

\begin{solution}
令 \(u=|a|>0\)，\(v=|b|>0\).因 \(ab>0\)，有 \(ab=uv\)，故原式等于 \(\dfrac{u^4+4v^4+1}{uv}\).由均值不等式，
\(u^4+4v^4+1=u^4+4v^4+\dfrac12+\dfrac12\ge4\sqrt[4]{u^4\cdot4v^4\cdot\dfrac12\cdot\dfrac12}=4uv\).所以原式不小于 \(4\).

当 \(u^4=4v^4=\dfrac12\) 时等号成立，例如取 \(a=2^{-1/4}\)，\(b=2^{-3/4}\)，满足 \(ab>0\).因此最小值为 \(4\).
\end{solution}

\begin{problem}
在 \(\triangle ABC\) 中，\(\angle A=60^\circ\)，\(AB=3\)，\(AC=2\). 若 \(\overrightarrow{BD}=2\overrightarrow{DC}\)，\(\overrightarrow{AE}=\lambda\overrightarrow{AC}-\overrightarrow{AB}\)（\(\lambda\in\R\)），且 \(\overrightarrow{AD}\cdot\overrightarrow{AE}=-4\)，则 \(\lambda\) 的值为\fillinblank{}.
\end{problem}

\begin{answer}
\(\dfrac{3}{11}\)
\end{answer}

\begin{solution}
令 \(\bs{u}=\overrightarrow{AB}\)，\(\bs{v}=\overrightarrow{AC}\).由已知，\(\bs{u}\cdot\bs{u}=9\)，\(\bs{v}\cdot\bs{v}=4\)，且 \(\bs{u}\cdot\bs{v}=3\cdot2\cos60^\circ=3\).

由 \(\overrightarrow{BD}=2\overrightarrow{DC}\)，得 \(D-B=2(C-D)\)，所以 \(3D=B+2C\)，从而 \(\overrightarrow{AD}=\dfrac13\bs{u}+\dfrac23\bs{v}\).又 \(\overrightarrow{AE}=\lambda\bs{v}-\bs{u}\)，故
\[
\overrightarrow{AD}\cdot\overrightarrow{AE}=\left(\dfrac13\bs{u}+\dfrac23\bs{v}\right)\cdot\left(\lambda\bs{v}-\bs{u}\right)=\dfrac{11\lambda-15}{3}
\]
令其等于 \(-4\)，得 \(\dfrac{11\lambda-15}{3}=-4\)，所以 \(\lambda=\dfrac3{11}\).
\end{solution}

\section{解答题}

\begin{problem}
在 \(\triangle ABC\) 中，内角 \(A\)，\(B\)，\(C\) 所对的边分别为 \(a\)，\(b\)，\(c\). 已知 \(a\sin A=4b\sin B\)，\(ac=\sqrt5\left(a^2-b^2-c^2\right)\).

\begin{enumerate}
\item 求 \(\cos A\) 的值；
\item 求 \(\sin(2B-A)\) 的值.
\end{enumerate}
\end{problem}

\begin{solution}
\begin{enumerate}
\item 由正弦定理，\(\dfrac{a}{\sin A}=\dfrac{b}{\sin B}\).因此 \(a\sin A=\dfrac{a^2}{2R}\)，\(b\sin B=\dfrac{b^2}{2R}\)，其中 \(R\) 为三角形外接圆半径.由题设得 \(a^2=4b^2\)，由于边长为正，故 \(a=2b\).

由余弦定理，\(a^2-b^2-c^2=-2bc\cos A\).代入第二个条件并除以 \(c>0\)，得 \(a=-2\sqrt5 b\cos A\).再代入 \(a=2b\)，得 \(\cos A=-\dfrac1{\sqrt5}=-\dfrac{\sqrt5}{5}\).

\item 因为 \(A\) 为钝角，所以 \(B\)，\(C\) 为锐角.由正弦定理及 \(a=2b\)，得 \(\sin B=\dfrac ba\sin A=\dfrac12\cdot\dfrac2{\sqrt5}=\dfrac1{\sqrt5}\)，从而 \(\cos B=\dfrac2{\sqrt5}\).于是 \(\sin2B=\dfrac45\)，\(\cos2B=\dfrac35\)，且 \(\sin A=\dfrac2{\sqrt5}\).

因此 \(\sin(2B-A)=\sin2B\cos A-\cos2B\sin A=\dfrac45\left(-\dfrac1{\sqrt5}\right)-\dfrac35\cdot\dfrac2{\sqrt5}=-\dfrac2{\sqrt5}=-\dfrac{2\sqrt5}{5}\).
\end{enumerate}
\end{solution}

\begin{problem}
电视台播放甲，乙两套连续剧，每次播放连续剧时，需要播放广告. 已知每次播放甲，乙两套连续剧时，连续剧播放时长，广告播放时长，收视人次如下表所示：

\begin{center}
\begin{tabular}{c|ccc}
 & 连续剧播放时长（分钟） & 广告播放时长（分钟） & 收视人次（万）\\
\hline
甲 & \(70\) & \(5\) & \(60\)\\
乙 & \(60\) & \(5\) & \(25\)
\end{tabular}
\end{center}

已知电视台每周安排的甲，乙连续剧的总播放时间不多于 \(600\) 分钟，广告的总播放时间不少于 \(30\) 分钟，且甲连续剧播放的次数不多于乙连续剧播放次数的 \(2\) 倍. 分别用 \(x\)，\(y\) 表示每周计划播出的甲，乙两套连续剧的次数.

\begin{enumerate}
\item 用 \(x\)，\(y\) 列出满足题目条件的数学关系式，并画出相应的平面区域；
\item 问电视台每周播出甲，乙两套连续剧各多少次，才能使总收视人次最多？
\end{enumerate}
\end{problem}

\begin{solution}
\begin{enumerate}
\item 由总播放时间、广告时间、次数关系及次数非负，得到
\[
\begin{cases}
7x+6y\le60,\\
x+y\ge6,\\
x-2y\le0,\\
x\ge0,\\
y\ge0.
\end{cases}
\]
其中第一式由 \(70x+60y\le600\) 同除以 \(10\) 得到，第二式由 \(5x+5y\ge30\) 同除以 \(5\) 得到.可行域是上述五个半平面的交集，边界依次为 \(x=0\) 上从 \((0,6)\) 到 \((0,10)\) 的线段、\(7x+6y=60\) 上从 \((0,10)\) 到 \((6,3)\) 的线段、\(x=2y\) 上从 \((6,3)\) 到 \((4,2)\) 的线段，以及 \(x+y=6\) 上从 \((4,2)\) 到 \((0,6)\) 的线段.因此可行域是顶点依次为 \((0,6)\)、\((0,10)\)、\((6,3)\)、\((4,2)\) 的封闭四边形，按这些边界即可画出相应区域.

\item 设总收视人次为 \(z\) 万，则目标函数为 \(z=60x+25y\).在可行域的四个顶点处分别计算：\(z(0,6)=150\)，\(z(0,10)=250\)，\(z(6,3)=435\)，\(z(4,2)=290\).最大值为 \(435\)，在 \((x,y)=(6,3)\) 处取得.因此每周应播出甲连续剧 \(6\) 次，乙连续剧 \(3\) 次.
\end{enumerate}
\end{solution}

\begin{problem}
如图，在四棱锥 \(P-ABCD\) 中，\(AD\perp\) 平面 \(PDC\)，\(AD\parallel BC\)，\(PD\perp PB\)，\(AD=1\)，\(BC=3\)，\(CD=4\)，\(PD=2\).

\begin{enumerate}
\item 求异面直线 \(AP\) 与 \(BC\) 所成角的余弦值；
\item 求证：\(PD\perp\) 平面 \(PBC\)；
\item 求直线 \(AB\) 与平面 \(PBC\) 所成角的正弦值.
\end{enumerate}

\bitmapfigure[max width=6.5cm]{img/2017/tianjin_liberal/q17_fig1.png}
\end{problem}

\begin{solution}
\begin{enumerate}
\item 因为 \(AD\parallel BC\)，所以异面直线 \(AP\) 与 \(BC\) 所成的角等于直线 \(AP\) 与 \(AD\) 所成角的锐角，即 \(\angle DAP\) 或其补角.又因为 \(AD\perp\) 平面 \(PDC\)，所以 \(AD\perp PD\).在直角三角形 \(PDA\) 中，\(AP=\sqrt{AD^2+PD^2}=\sqrt5\)，故所成角的余弦值为 \(\dfrac{AD}{AP}=\dfrac1{\sqrt5}=\dfrac{\sqrt5}{5}\).

\item 由 \(AD\perp\) 平面 \(PDC\) 且 \(PD\) 在平面 \(PDC\) 内，得 \(AD\perp PD\).又由 \(AD\parallel BC\)，得 \(PD\perp BC\).已知 \(PD\perp PB\)，而 \(PB\) 与 \(BC\) 是平面 \(PBC\) 内相交的两条直线，所以 \(PD\perp\) 平面 \(PBC\).

\item 过点 \(D\) 作 \(DF\parallel AB\)，交 \(BC\) 于点 \(F\).由 \(AD\parallel BC\) 得 \(AD\parallel BF\)，所以四边形 \(ABFD\) 是平行四边形，进而 \(DF=AB\)，且 \(BF=AD=1\).因此 \(CF=BC-BF=2\).

由 \(AD\perp\) 平面 \(PDC\) 得 \(AD\perp DC\)，再由 \(AD\parallel BC\) 得 \(BC\perp DC\)，所以三角形 \(DCF\) 在 \(C\) 处直角，\(DF=\sqrt{DC^2+CF^2}=\sqrt{4^2+2^2}=2\sqrt5\).

由第二问，\(PD\perp\) 平面 \(PBC\)，故 \(PF\) 是 \(DF\) 在平面 \(PBC\) 上的射影，直线 \(DF\) 与平面 \(PBC\) 所成的角为 \(\angle DFP\).由于 \(DF\parallel AB\)，这也就是直线 \(AB\) 与平面 \(PBC\) 所成的角.三角形 \(DPF\) 在 \(P\) 处直角，故其正弦为 \(\dfrac{PD}{DF}=\dfrac2{2\sqrt5}=\dfrac{\sqrt5}{5}\).
\end{enumerate}
\end{solution}

\begin{problem}
已知 \(\{a_n\}\) 为等差数列，前 \(n\) 项和为 \(S_n\)（\(n\in\N^*\)），\(\{b_n\}\) 是首项为 \(2\) 的等比数列，且公比大于 \(0\)，\(b_2+b_3=12\)，\(b_3=a_4-2a_1\)，\(S_{11}=11b_4\).

\begin{enumerate}
\item 求 \(\{a_n\}\) 和 \(\{b_n\}\) 的通项公式；
\item 求数列 \(\{a_{2n}b_n\}\) 的前 \(n\) 项和（\(n\in\N^*\)）.
\end{enumerate}
\end{problem}

\begin{solution}
\begin{enumerate}
\item 设等差数列 \(\{a_n\}\) 的公差为 \(d\)，等比数列 \(\{b_n\}\) 的公比为 \(q\).由 \(b_1=2\)，有 \(b_2=2q\)，\(b_3=2q^2\)，所以 \(2q+2q^2=12\)，即 \(q^2+q-6=0\).因 \(q>0\)，得 \(q=2\)，故 \(b_n=2\cdot2^{n-1}=2^n\).

由 \(b_3=a_4-2a_1\) 得 \(8=(a_1+3d)-2a_1=3d-a_1\).又 \(S_{11}=\dfrac{11}{2}(2a_1+10d)=11(a_1+5d)\)，而 \(11b_4=11\cdot16\)，所以 \(a_1+5d=16\).解方程组
\(3d-a_1=8\)，\(a_1+5d=16\)，得 \(a_1=1\)，\(d=3\).因此 \(a_n=1+3(n-1)=3n-2\).

\item 设所求前 \(n\) 项和为 \(T_n\).由上问，\(a_{2k}=6k-2\)，\(b_k=2^k\)，故
\(T_n=\sum_{k=1}^n(6k-2)2^k=6\sum_{k=1}^n k2^k-2\sum_{k=1}^n2^k\).其中 \(\sum_{k=1}^n2^k=2^{n+1}-2\).令 \(U_n=\sum_{k=1}^n k2^k\)，将 \(2U_n=\sum_{k=2}^{n+1}(k-1)2^k\) 与 \(U_n\) 错位相减，得
\[
2U_n-U_n=n2^{n+1}-2-\sum_{k=2}^{n}2^k=n2^{n+1}-2-\left(2^{n+1}-4\right)=(n-1)2^{n+1}+2
\]
所以
\[
T_n=6\left((n-1)2^{n+1}+2\right)-2\left(2^{n+1}-2\right)=(3n-4)2^{n+2}+16
\]
\end{enumerate}
\end{solution}

\begin{problem}
设 \(a\)，\(b\in\R\)，\(|a|\le 1\). 已知函数 \(f(x)=x^3-6x^2-3a(a-4)x+b\)，\(g(x)=\e^x f(x)\).

\begin{enumerate}
\item 求 \(f(x)\) 的单调区间；
\item 已知函数 \(y=g(x)\) 和 \(y=\e^x\) 的图象在公共点 \((x_0,y_0)\) 处有相同的切线.
\begin{enumerate}
\item 求证：\(f(x)\) 在 \(x=x_0\) 处的导数等于 \(0\)；
\item 若关于 \(x\) 的不等式 \(g(x)\le \e^x\) 在区间 \([x_0-1,x_0+1]\) 上恒成立，求 \(b\) 的取值范围.
\end{enumerate}
\end{enumerate}
\end{problem}

\begin{solution}
\begin{enumerate}
\item \(f'(x)=3x^2-12x-3a(a-4)=3(x-a)(x-4+a)\).由 \(|a|\le1\) 得 \(a<4-a\).因此当 \(x<a\) 或 \(x>4-a\) 时 \(f'(x)>0\)，当 \(a<x<4-a\) 时 \(f'(x)<0\).所以 \(f(x)\) 的单调递增区间为 \((-\infty,a)\)、\((4-a,+\infty)\)，单调递减区间为 \((a,4-a)\).

\item
\begin{enumerate}
\item 公共点条件给出 \(g(x_0)=\e^{x_0}\)，即 \(\e^{x_0}f(x_0)=\e^{x_0}\)，所以 \(f(x_0)=1\).又 \(g'(x)=\e^x(f(x)+f'(x))\)，而两图象在该点的切线相同，故 \(g'(x_0)=\e^{x_0}\)，从而 \(f(x_0)+f'(x_0)=1\).结合 \(f(x_0)=1\)，得到 \(f'(x_0)=0\).

\item 因为 \(\e^x>0\)，不等式 \(g(x)\le\e^x\) 等价于 \(f(x)\le1\).由上问，\(f(x_0)=1\)，\(f'(x_0)=0\)，所以 \(x_0=a\) 或 \(x_0=4-a\).若 \(x_0=4-a\)，则在其右侧足够小的邻域内 \(f\) 递增，因而有 \(f(x)>f(x_0)=1\)，与区间上的恒成立条件矛盾，故 \(x_0=a\).

由于 \(|a|\le1\)，有 \(a+1<4-a\).由第一问，\(f\) 在 \([a-1,a]\) 上递增，在 \([a,a+1]\) 上递减，所以 \(f(x)\le f(a)=1\) 对 \(x\in[a-1,a+1]\) 恒成立.于是条件等价于 \(f(a)=1\).

计算得 \(f(a)=a^3-6a^2-3a^2(a-4)+b=-2a^3+6a^2+b\)，所以 \(b=1+2a^3-6a^2\)，其中 \(-1\le a\le1\).令 \(t(a)=1+2a^3-6a^2\)，则 \(t'(a)=6a(a-2)\).在 \([-1,0]\) 上 \(t'(a)>0\)，在 \([0,1]\) 上 \(t'(a)<0\)，所以最大值为 \(t(0)=1\)，最小值为 \(t(-1)=-7\)（另有 \(t(1)=-3\)）.因此 \(b\in[-7,1]\).
\end{enumerate}
\end{enumerate}
\end{solution}

\begin{problem}
已知椭圆 \(\dfrac{x^2}{a^2}+\dfrac{y^2}{b^2}=1\)（\(a>b>0\)）的左焦点为 \(F(-c,0)\)，右顶点为 \(A\)，点 \(E\) 的坐标为 \((0,c)\)，\(\triangle EFA\) 的面积为 \(\dfrac{b^2}{2}\).

\begin{enumerate}
\item 求椭圆的离心率；
\item 设点 \(Q\) 在线段 \(AE\) 上，\(|FQ|=\dfrac32c\)，延长线段 \(FQ\) 与椭圆交于点 \(P\)，点 \(M\)，\(N\) 在 \(x\) 轴上，\(PM\parallel QN\)，且直线 \(PM\) 与直线 \(QN\) 间的距离为 \(c\)，四边形 \(PQNM\) 的面积为 \(3c\).
\begin{enumerate}
\item 求直线 \(FP\) 的斜率；
\item 求椭圆的方程.
\end{enumerate}
\end{enumerate}
\end{problem}

\begin{solution}
\begin{enumerate}
\item 由 \(A=(a,0)\)、\(F=(-c,0)\)、\(E=(0,c)\)，得 \(|FA|=a+c\)，点 \(E\) 到 \(x\) 轴的距离为 \(c\)，所以
\(\dfrac12(a+c)c=\dfrac{b^2}{2}\)，即 \((a+c)c=b^2\).又 \(b^2=a^2-c^2=(a+c)(a-c)\)，故 \(c=a-c\)，即 \(a=2c\).于是 \(b^2=3c^2\)，椭圆的离心率为 \(e=\dfrac ca=\dfrac12\).

\item
\begin{enumerate}
\item 由上问，\(A=(2c,0)\)，\(E=(0,c)\)，所以直线 \(AE\) 的方程为 \(x+2y=2c\).设直线 \(FP\) 的方程为 \(x=my-c\)，其中 \(m>0\)，则其斜率为 \(\dfrac1m\).与直线 \(AE\) 联立，得
\(Q\left(\dfrac{2c(m-1)}{m+2},\dfrac{3c}{m+2}\right)\).因此
\(FQ=\dfrac{3c\sqrt{m^2+1}}{m+2}=\dfrac{3c}{2}\)，即 \(2\sqrt{m^2+1}=m+2\).平方并整理得 \(3m^2-4m=0\).\(m=0\) 时交点不在线段 \(AE\) 上，故 \(m=\dfrac43\)，从而直线 \(FP\) 的斜率为 \(\dfrac34\).

\item 由 \(m=\dfrac43\)，直线 \(FP\) 为 \(3x-4y+3c=0\).又 \(a=2c\)，\(b^2=3c^2\)，椭圆方程为 \(\dfrac{x^2}{4c^2}+\dfrac{y^2}{3c^2}=1\).联立直线方程，消去 \(y\) 得
\(7x^2+6cx-13c^2=0\)，所以 \(x=-\dfrac{13c}{7}\) 或 \(x=c\).延长线段 \(FQ\) 与椭圆交于 \(P\)，取在射线 \(FQ\) 上的交点，故 \(P\left(c,\dfrac32c\right)\).由此 \(FP=\dfrac52c\)，而 \(FQ=\dfrac32c\)，所以 \(PQ=c\).

由题设，两平行直线 \(PM\)、\(QN\) 的距离为 \(c\).线段 \(PQ\) 的两个端点分别在这两条直线上，且 \(|PQ|=c\)，故 \(PQ\) 正好是两平行直线的公垂线，即 \(PM\perp FP\)、\(QN\perp FP\).设 \(\theta=\angle QFN\)，则 \(\tan\theta=\dfrac34\).在直角三角形 \(FQN\) 中，
\(QN=FQ\tan\theta=\dfrac{3c}{2}\cdot\dfrac34=\dfrac{9c}{8}\)，故 \(S_{\triangle FQN}=\dfrac12\cdot\dfrac{3c}{2}\cdot\dfrac{9c}{8}=\dfrac{27c^2}{32}\).在直角三角形 \(FPM\) 中，\(\tan\angle PFM=\dfrac34\)，所以 \(PM=FP\tan\angle PFM=\dfrac{5c}{2}\cdot\dfrac34=\dfrac{15c}{8}\)，故 \(S_{\triangle FPM}=\dfrac{75c^2}{32}\).

四边形 \(PQNM\) 的面积为两三角形面积之差，所以 \(\dfrac{75c^2}{32}-\dfrac{27c^2}{32}=3c\).由 \(c>0\) 得 \(c=2\).于是 \(a=4\)，\(b^2=12\)，椭圆方程为 \(\dfrac{x^2}{16}+\dfrac{y^2}{12}=1\).
\end{enumerate}
\end{enumerate}
\end{solution}
